% TODO(PROSE): Final abstract not written in M6.
% Structured skeleton only.

\begin{abstract}
\textbf{Problem:} CPU-side PMP and DMA-side IOPMP-style mechanisms protect different transaction origins in RISC-V SoCs.

\textbf{Research question:} Whether local access-control correctness composes into an end-to-end protected-memory isolation guarantee.

\textbf{Method:} RTL simulation, randomized tests, bounded model checking, and induction/PDR where applicable, with assumption sensitivity.

\textbf{Verified result:} Selected normal-operation CPU/DMA access-control properties (e.g., SP-02, SP-04) are formally established under documented assumptions.

\textbf{Compositional result:} Under RST-A fail-open initialization, an unsafe early-DMA protected-memory state before \texttt{secure\_ready} is reachable (RESET\_ASSUMPTION\_DEPENDENCY). Under RST-B fail-closed defaults, the corresponding state is unreachable (FORMAL\_PROOF via k-induction).

\textbf{Conclusion:} Compositional isolation depends on initialization, reset ordering, requester identity, and transaction-admission invariants---not local rules alone.

% EVIDENCE: results/m510/m510_summary.md
\end{abstract}
